% SEGMENT OLP-0585-S01
\documentclass[../../../include/open-logic-section]{subfiles}

\begin{document}

\olfileid{sth}{cardinals}{hp}
\olsection{பிற்சேர்க்கை: ஹியூமின் கொள்கை}

\olref[cp]{sec} இல் காண்டோரின் கொள்கையைக் கூறினோம். அது:
\begin{align*}
	\card{A} = \card{B} & \text{ iff } A \approx B.
\intertext{இது இப்போது \emph{ஹியூமின் கொள்கை} எனப்படும் பின்வரும்
கூற்றை மிகவும் ஒத்திருக்கிறது:}
	\fregenum{x} {F(x)} = \fregenum{x}{G(x)} & \text{ iff } F \sim G
\end{align*}
இங்கு `$F \sim G$' என்பதன் சுருக்கப் பொருள், $F$s எத்தனை உள்ளனவோ
அதே எண்ணிக்கையில் $G$s உம் உள்ளன என்பதே. அதாவது, $F$s க்கும்
$G$s க்கும் இடையில் ஓர் இருபுறச் சார்பை அமைக்க முடியும்; முறையாக:
\begin{align*}
	\exists R(&\forall v\forall y(Rvy \lif (Fv \land Gy)) \land {}\\
		&\forall v(Fv \lif \lexists![y][Rvy]) \land {}\\
		&\forall y(Gy \lif \lexists![v][Rvy]))
\end{align*}

% SEGMENT OLP-0585-S02
ஆனால் ஹியூமின் கொள்கைக்கும் காண்டோரின் கொள்கைக்கும் வகை வேறுபாடு
உள்ளது. காண்டோரின் கொள்கையில் ``$A$'', ``$B$'' என்பவை
\emph{கணங்களைக்} குறிக்கும் முதல்-வரிசைப் பதங்கள்.
ஹியூமின் கொள்கையில் ``$F$'', ``$G$'', ``$R$'' என்பவை
முதல்-வரிசைப் பதங்கள் \emph{அல்ல}; அவை \emph{பயனிலை
நிலைகளில்} வருகின்றன. (ஒருவேளை அவை \emph{பண்புகளைக்}
குறிக்கலாம்.) எனவே ஹியூமின் கொள்கையை இவ்வாறு விளக்கலாம்:
$F$s இன் எண் $G$s இன் எண்ணுக்குச் சமம் எனும்போதும் அப்போதுதான்
$F$s மற்றும் $G$s இடையே ஓர் இருபுறச் சார்பு உள்ளது.
இது \emph{ஹியூமின் கொள்கை} எனப்படுகிறது; ஏனெனில் ஹியூம்
ஒருமுறை இவ்வாறு எழுதினார்:
\begin{quote}
  இரண்டு எண்களை இவ்வாறு இணைத்து, ஒன்றின் ஒவ்வோர் அலகுக்கும்
  மற்றொன்றின் ஓர் அலகு எப்போதும் பொருந்தினால், அவை சமம் என்று
  நாம் கூறுகிறோம்.
  \citep[Pt.III Bk.1 \S1]{Hume1740}
\end{quote}
ஹியூமின் இந்தப் பகுதியை மேற்கோள் காட்டிய \citet[\S63]{Frege1884},
பின்னர் (நாம் அழைக்கும்) ஹியூமின் கொள்கையைப் பெருமளவில்
வரைந்து காட்டினார். இதன் மூலம் இக்கொள்கை இன்றைய கணிதத்
தருக்கவியலில் முக்கியத்துவம் பெற்றது.

% SEGMENT OLP-0585-S03
ஹியூமின் கொள்கைக்கும் அடிப்படை விதி~V க்கும் உள்ள கட்டமைப்புச்
சார்ந்த ஒற்றுமையைக் கவனிக்கவும். \olref[story][blv]{sec} இல்
அதைப் பின்வருமாறு எழுதினோம்:
\[
	\fregeext{x}{F(x)} = \fregeext{x}{G(x)} \text{iff } \lforall[x][(F(x) \liff G(x))].
\]
இப்போது இதைப் பற்றிச் சில விளக்கங்களும் ஒப்பீடும் உதவும்.

ஹியூமின் கொள்கை அல்லது அடிப்படை விதி~V போன்ற ஒரு கொள்கையை
\emph{முன்னறிவுறு} முறையிலோ \emph{முன்னறிவுறாத} முறையிலோ
புரிந்துகொள்ளலாம் (\olref[story][predicative]{sec} ஐ நினைவுகூர்க).
அடிப்படை விதி~V ஐ முன்னறிவுறாத முறையில் வாசித்தால், ஒவ்வொரு~$F$
க்கும் $\fregeext{x}{F(x)}$ என்ற பொருள், அந்த விதியை வகுக்கும்போது
நாம் பயன்படுத்திய அளவையிடல் களத்திலேயே அடங்கும்.
அதேபோல, ஹியூமின் கொள்கையை முன்னறிவுறாத முறையில் வாசித்தால்,
ஒவ்வொரு~$F$ க்கும் $\fregenum{x}{F(x)}$ என்ற பொருள்,
அக்கொள்கையை வகுக்கும்போது பயன்படுத்திய அளவையிடல் களத்திலேயே
அடங்கும். மாறாக, \emph{முன்னறிவுறு} புரிதலில்,
$\fregeext{x}{F(x)}$ மற்றும்~$\fregenum{x}{F(x)}$ ஆகிய
பொருள்கள் \emph{வேறொரு} களத்தைச் சேர்ந்திருக்கும்.

% SEGMENT OLP-0585-S04
அடிப்படை விதி~V ஐ முன்னறிவுறாத முறையில் வாசித்தால்,
முறைசாராப் பண்புவழிக் கணமாக்கலின் வழியாக அது முரண்பாட்டுக்கு
இட்டுச் செல்கிறது (விவரங்களுக்கு \olref[story][blv]{sec} ஐப்
பார்க்கவும்). முறைசாராப் பண்புவழிக் கணமாக்கலைப் போலவே,
அதனை \emph{முன்னறிவுறு} முறையில் வாசித்தால் முரண்பாடற்றதாக
ஆக்கலாம். ஆனால் அப்போது அதிலிருந்து நாம் விரும்பிய எல்லா
விளைவுகளும் கிடைக்காமல் போகலாம்.

ஹியூமின் கொள்கையையோ \emph{முரண்பாடின்றி} முன்னறிவுறாத முறையில்
வாசிக்க முடியும். அவ்வாறு வாசிக்கும்போது அது கணிசமான ஆற்றல்
கொண்டது.

% SEGMENT OLP-0585-S05
இதற்கு எடுத்துக்காட்டாக, எந்தப் பொருளும் நிறைவேற்றாத ``$x \neq x$''
என்ற பயனிலையைக் கருதுங்கள். ஹியூமின் கொள்கை இப்போது
$\# x( x\neq x)$ என்ற பொருளைத் தருகிறது. இதை $0$ என்ற எண்ணாகக்
கொள்ளலாம். \emph{முன்னறிவுறாத} புரிதலில் --- அதில்
\emph{மட்டுமே} --- இந்தப் பொருள் $0$ நமது முதன்மை அளவையிடல்
களத்திலேயே அடங்கும். ஆகவே ``$x = 0$'' என்ற பயனிலைக்குப்
ஹியூமின் கொள்கையைப் பயன்படுத்தி $\#x (x = 0)$ என்ற பொருளைப்
பெறலாம். இதை $1$ என்ற எண்ணாகக் கொள்ளலாம். மேலும் $0 \neq 1$
என்று ஹியூமின் கொள்கை தருகிறது; ஏனெனில் தன்னுடன் சமமில்லாத
பொருள்களுக்கும் $0$ உடன் சமமான பொருள்களுக்கும் இடையில் இருபுறச்
சார்பு இருக்க முடியாது. முதலாவது வகையில் ஒன்றுமில்லை; இரண்டாவதில்
ஒரே ஒரு பொருள் உள்ளது. மீண்டும் முன்னறிவுறாத முறையில்
செயல்படும்போது, $1$ உம் நமது முதன்மை அளவையிடல் களத்தில்
அடங்கும். எனவே ``$(x = 0 \lor x = 1)$'' என்ற பயனிலைக்குப்
ஹியூமின் கொள்கையைப் பயன்படுத்தி $\#x(x = 0 \lor x = 1)$
என்ற பொருளைப் பெறலாம். இதை $2$ என்ற எண்ணாகக் கொள்ளலாம்;
$0\neq 2$ மற்றும் $1 \neq 2$ எனத் தொடர்ந்து காட்டலாம்.

% SEGMENT OLP-0585-S06
ஆகவே, முன்னறிவுறாத முறையில் எடுத்துக்கொண்டால்,
\emph{முடிவுறாத எண்ணிக்கையில் பொருள்கள் உள்ளன} என்று
ஹியூமின் கொள்கை தருகிறது. இதனால் \emph{புதுப் ஃபிரேகிய
தருக்கவாதிகள்} இக்கொள்கையை எண்கணிதத்தின் அடித்தளமாகக்
கொள்ளத் தூண்டப்பட்டுள்ளனர்.

ஆனால் ஃபிரேகே \emph{தாமே} ஹியூமின் கொள்கையை எண்கணிதத்தின்
அடித்தளமாக எடுத்துக்கொள்ளவில்லை. மாறாக, ஒரு வெளிப்படையான
வரையறையிலிருந்து அதை நிறுவினார்: $\fregenum{x}{F(x)}$ என்பது
$F \sim \Phi$ என்ற கருத்தின் விரிவாக வரையறுக்கப்பட்டது.
இன்றைய மொழியில் இதனை $\fregenum{x}{F(x)} = \Setabs{G}{F \sim G}$
என்று எழுத முயலலாம்; ஆனால் அப்போது முறைசாராப் பண்புவழிக்
கணமாக்கலின் சிக்கல்களுக்கே திரும்பி வருவோம்.

\end{document}
